\section{The Conservatives of Reflection}

At one time, there was no such political movement similar to the ``right'', because what is now called the ``right'', properly understood, was the natural organization of traditional societies. It was only the spirit of rebellion, which has marked all recent Western history, that created this dichotomy between the so-called left and right. Note that by this standard there is really nothing of the ``right'' remaining in the West and it should not be confused with contemporary ``conservative'' movements.

Historically, the terms derived from the two opposing forces resulting from the French revolution. That led to two distinct, but related, currents of conservative thought. One, which we have focused on, derives from the continental stream beginning with \textbf{Joseph de Maistre} and continuing with \textbf{Bonald}, \textbf{Louis Veuillot}, \textbf{Donoso Cortes}, and many similar critics of liberalism and revolution.

The other current, almost exclusively in English, derives from \textbf{Edmund Burke}; this current has been amply documented in the works of \textbf{Russell Kirk}, particularly in \emph{The Conservative Mind}. This current is neglected and nearly forgotten; it certainly has no voice in the public square. Instead, in recent times, new currents have arisen, with little or no connection to these two primary currents. Recognizing this, they are ugly prefixes like ``neo'', ``new'', or ``alternative''. Such terms give me agita\footnote{\url{https://blogs.law.harvard.edu/ethicalesq/what-is-agita/}} since the prefix simply sounds like ``non'' to my mind. If you are heading in the right direction, then any innovations or novelties must necessarily be in another direction. If not, they are superfluous.

Now I can understand that changing historical circumstances may require new presentations, given that the past cannot repeat itself. However, I believe that is not what is going on, since these novelties seem to embrace the mentality of the modern world that they claim to oppose. Far from being anti-intellectual, the right is the true intellect. Hence, in a schematic way with minimal explanation, I will list its ways of thinking and understanding. Perhaps someone can demonstrate that these novelties accept all these principle, and thus are not novelties at all. Otherwise, it will be helpful to point out their deviations. It is important to stay on the level of principles rather than specific issues; for to say that a ``conservative'' issue is this or that is often wrong unless there is a reason behind it.

\paragraph{Three principles of Knowledge}


\begin{itemize}
\item \textbf{Transcendence}: conformance of the mind to transcendent ideas
\item \textbf{Reality}: conformance of the mind to the world
\item \textbf{Integrity}: conformance of the interior man to the exterior man in his life, speech, and actions
\end{itemize}


The modern mind rejects at least the first two principles, and probably the third. For it, there are no transcendent ideas in the Mind of God; rather ideas are concepts formed in man's mind. Since Kant, the external world is held to be unknowable in itself; instead, we form ideas about it either through science or ideologies.

\paragraph{Order and Chaos}


There are two ways to relate these notions. The first is traditional, the second is modern, and is the logical consequence of scientism and atheism.

\begin{itemize}
\item Order, or Logos, is the essential aspect, chaos is the accident
\item Chaos is the essential aspect, and order is an accident
\end{itemize}


\paragraph{Ideology: the Opposite of Knowledge}


In rejecting the knowledge of reality, the modern mind is prone to ideology. That is, since there is nothing real that it must conform to, it is free to create grandiose schemes. There is really no common ground for discussion between the two worldviews. It seems to me that all the ``new'' rights are ideological. For example, there are the incessant debates about religion, as though one could pick and choose a religion based on how well it conforms to their ideology. Instead, political thought must follow spiritual reality, not the other way around.

\paragraph{Conservative Premises}


In the Introduction to The Portable Conservative Reader, Russell Kirk expounds on six conservative premises that he derives from Edmund Burke. These are not necessarily metaphysical principles but rather prudential principles intended to guide practical politics. Hence, their applicability will vary for different circumstances.

\begin{itemize}


\item \textbf{Transcendent Moral Order}

There is always a ``lawgiver'', with divine sanction in a traditional society. In the West, that has taken the form of natural law.

\item \textbf{Social continuity}

The order of society is the product of a ``long and painful social experience, the results of centuries of trial and reflection and sacrifice.'' This is not to say that change is not necessary, but it needs to be gradual.

\item \textbf{Prescription}

This is adherence to the ``wisdom of our ancestors''. The first tactic of the revolutionary is always to attack the ancestors. The new rightists, in their constant attack on the foundations of Western civilization, reveal themselves as revolutionaries.

\item \textbf{Prudence}

Plato and Burke agree that prudence is the chief virtue of the statesman. The long run consequences of a policy must be taken into account, rather than temporary advantage or popularity.

\item \textbf{Variety}

Conservatives ``feel affection for the proliferating intricacy of long-established social institutions and modes of life as distinguished from the narrowing uniformity and deadening egalitarianism of radical systems.'' True and healthy diversity rests of many sorts of inequality.

\item \textbf{Imperfectibility}

``Human nature suffers irremediably from certain faults.'' This is the fundamental axiom of the science of man. The corollary is that ``no perfect social order ever can be created.'' All schemes that depend on men acting intelligently and with good will in all circumstances will fail. A fortiori, any system that requires the appearance of a ``new man'' will fail miserably.
\end{itemize}


\paragraph{Three Types of Conservatives}


\textbf{Walter Bagehot} distinguishes three types of conservatives.


\begin{itemize}
\item The Conservatives of Enjoyment
\item The Conservatives of Order
\item The Conservatives of Reflection
\end{itemize}


The \textbf{conservatives of enjoyment} like the ``old ways'' for their own sake. Hence, they prefer ancestral customs in manners, dress, religion, thought, and so on. For this they need no intellectual conviction to justify that preference. Much is taken on ``animal faith'', supported by established and ancient traditions. Their primary motivation is loyalty.

The \textbf{conservatives of order} are the petite bourgeoisie who want to hold onto their property, land, family, etc. Thus they support policies that will protect their holdings and wealth, as well as ``law and order''. They are often mocked, but yet they form the bedrock of a society. Their primary motivation is fear.

The \textbf{conservatives of reflection} know the reasons for loyalty, order, etc., in short, the principles mentioned above. In Bagehot's time, he claimed there were too few of this type and little has changed since then. The first two types of conservatives are on the soapbox and their lack of intellectual sophistication makes them easy targets for the forces of revolution.

So what have the novelties brought us? They reject the ``old ways'' and are not loyal to our ancestors, or else they count as ancestors only those from some remote, unknowable, and imaginary antiquity. In actuality, they prefer the ways, mores, and beliefs of the modern world and have often admitted it. The task of the new conservatives of reflection is to recover those ideas and restate them for modern ears.

\paragraph{Addendum 18 August 2022}


This may be of interest to so called conservatives in the USA. The intellectual level of the American right is embarrassingly deficient; no wonder it is called the ``stupid party''. Back in the 70s and 80s, the right was based on more serious thinking. \textbf{John East}, in The American Conservative Movement: The Philosophical Foundations, identified seven thinkers, now mostly forgotten or ignored. Although they don't all fit together nicely, reading them indicates a serious commitment. This is the list, without comment.

\begin{enumerate}
\item Russell Kirk
\item Richard Weaver
\item Frank S. Meyer
\item Willmoore Kendall
\item Leo Strauss
\item Eric Voegelin
\item Ludwig von Mises
\end{enumerate}


We would also add James Burnham to this list.


\flrightit{Posted on 2013-08-18 by Cologero}

\begin{center}* * *\end{center}

\begin{footnotesize}
\begin{sffamily}
\texttt{Avery Morrow on 2013-08-18 at 22:53 said:}

Where is the conservatism of Don Fabrizio?

\hfill

\texttt{Cologero on 2013-08-18 at 23:04 said:}

The first type, of course. See \textit{NTELLECTUAL CONSERVATISM}\footnote{\url{http://oll.libertyfund.org/?option=com\_staticxt&staticfile=show.php\%3Ftitle=2167&chapter=200949&layout=html&Itemid=27}}.

\hfill

\texttt{Cassiodorus on 2013-08-19 at 00:26 said:}

Many contemporary rightists argue that there is, in fact, a genuine conservatism that isn't \emph{entirely} rooted in the pre-modern world, but flows larglely from the Scottish Enlightenment- not the continental strain that burst out of France. Writers like Russell Kirk defend the view that this tradition of republicanism and ``ordered liberty'' is largely the working out of conservative principles. The most basic claim that modern rightists never fail to communicate is the idea that liberty is central to conservatism because it rests on the theological principle of \emph{imago Dei}. Does this ``fusion'' of tradition and liberty really hold up?

\hfill

\texttt{Cologero on 2013-08-19 at 04:20 said:}

I agree, the notion of liberty is central; we see this particularly in Evola and we don't need to repeat everything here. Of course, we need to have the proper understanding of ``liberty'', not the modern one. As you point out, it is based on the Imago Dei, whereas the modern notion is that man replaces God.

\hfill

\texttt{Avery Morrow on 2013-08-19 at 21:39 said:}

@verbavolant: Polytheistic religion, as revived today, is a force of anti-tradition. You say that you enjoyed the music of the fairies. Which kind? A homosexual will gleefully call himself a fairy, but he will never gain the courage to earnestly confess his sins to a priest. Do read \emph{My Life Among the Deathworks} to learn more about the relationship between the classical and modern worlds.

\hfill

\texttt{August on 2013-08-20 at 01:58 said:}

All this airy harping upon færies, deserts and ancient woods might make one thirst for a pint of that famed Irish stout. Is this what was meant by raw food?

\hfill

\texttt{Logres on 2013-08-20 at 07:53 said:}

I'll second the Phillip Rieff recommendation! That book changed my mind in a way few others have. Verbolant, you might check this out:

\url{http://www.roadtoemmaus.net/back\_issue\_articles/RTE\_15/Brittanys\_Celtic\_Past.pdf}

I have a lot of sympathy for where you come from, as I started out in a ``Celtic mood'' myself! And I almost gave up Christianity, at several points, in despair, not understanding that Christianity actually incorporated the polytheism (the best of it) inside of it: if you want to be polytheistic, go back to the saints. It died off, mostly, outside the Church. Sure, I suppose remnants exist, but the stream went underground inside the Church. If you really want to ``get back to it'', the best path might be to pray to a saint (perhaps you already do), \& follow the Spirit. The revenge that the resurgent polytheists want to accomplish in restoring the ``icons'' is best \& only done through a path of forgiveness -- the Church would be transformed, were the best of ``pagans'' to enter it with full awareness of their heritage, and full humility of what else there is to learn. It's the balance of those two that will restore the West, nothing else. Strange times we live in.

\hfill

\texttt{Saladin on 2013-08-20 at 11:56 said:}

``The Conservatives of Enjoyment''

That is what previously was called Nobility!

\hfill

\texttt{Jacob on 2013-08-20 at 13:22 said:}

I have been reading Kirk's book as well. I thought his explanations on reason and tradition were excellent. With all the New Atheists running around claiming they have dispelled religious belief by ``reason'', it's interesting to known that reason can not even exist independent of Tradition. Worth the entire book IMO.

\hfill

\texttt{Ash on 2013-08-20 at 22:56 said:}

One of the most useful things about the old Right for us is the fact that they were writing in a time when many aspects of the Traditional order were still accepted as values and practiced in society. Though it was in decline and serious need of reform and restoration, it allows us to see what these principles look like when they are followed in the world and in daily life. Following that line of thought, the writings of someone like Thomas Merton may allow us to see how these values can be followed ``amidst the ruins''. What is left it to remain standing and begin to slowly rebuild, first within ourselves and then the world around us. The conservative of reflection must apply what has been learned.

h.ontologia, the author you mentioned definitely looks worth pursuing. Do you have any idea if he ever talked about socialism's eventual collapse? In addition to Merton, I highly recommend Alan Watts as an author. Although something of a mixed bag, his thoughts on how the principle of quality over quantity can manifest in simple daily phenomena like cooking or clothing are thought-provoking.

\hfill

\texttt{Cassiodorus on 2013-08-22 at 20:12 said:}

I have to pre-apologize for this derailment and because I've beat this horse before. But I'm having a tough time getting away from this both spiritually and ``intellectually''. I often read blogs that focus on Aristotelian Thomism- I recently made a comment that the ``spiritual level is known, not through empirical evidence or a process of reasoning, but rather through a direct knowing called Intuition or gnosis''- Cologero

A Thomist commenter replied with what seemed to be contempt for the notion of supra-rational knowledge. \emph{``Aristotle and Thomas both say that knowing comes from intellection upon sensory data as source material. There is, for humans, no knowing that is prior to sensation, and there is no adequate basis for a claim that we DO have prior knowledge- nobody can ask of a fetus pre-sensory ``tells us what you know.''}

As I understand it, intellection and gnosis are essentially the same, or at least, related terms. What's more is that the Angelic Doctor had a mystical experience towards the end of his life that he said transcended his dialectic. There are several things that are bouncing around for me- Revelation, faith, reason, and gnosis. I've been told by many a faithful Catholic that the Perennialist view that pure metaphysics transcends doctrine from within is exactly what the gnostic heresy is.

Cologero, go easy on me.

\hfill

\texttt{Cologero on 2013-08-22 at 20:41 said:}

Yes, Cassiodorus, at least you understand what is at stake. Do you find anything attractive or compelling about what that commentator told you? Who said anything about prior knowledge? Knowing comes from ``seeing'' that the sensory data is actually the reflection of a divine idea. As you point out, ``mystical'' experiences cannot be ``told'', obviously because they transcend the discursive mind.

In everyday life, what does it mean to ``get a joke'' or to suddenly understand a math problem? You can look at it (sensory data), discuss it (discursive mind), but the act of ``getting it'' is something direct, unmistakeable, infallible, and not like perception or thinking.

More basically, what exactly constitutes ``sensory data''? Do we really need all that high powered metaphysical apparatus to recognize a rabbit in the glen? Sensory data is much more complex and involves a hierarchy of higher level systems. How do we understand human action? I wrote a post here, but I can't recall the title, about the sensory data of one man holding a gun pointed at another. Is that an illegal robbery, a legal arrest, or actors filming the movie? The sensory data is certainly insufficient to answer that question.

Keep in mind that the Medievals also recognized several inner senses, in addition to the five outer senses. A thought is sensory data, an image, a dream, etc. How do you understand a dream? Can't it suddenly reveal itself to you? If so, this is beyond philosophy.

How do we understand human events? There are great movements of thought and it is the \textit{Spirit of Prophecy}\footnote{\url{http://www.meditationsonthetarot.com/the-future-of-prophecy}} that can reveal their inner meaning. Some may reduce such understanding to mere ``conspiracy theories''. The point is that the sensory data alone tells us nothing; the spirit of prophecy as St Pope Gregory explains it, transcends sensory data and discursive knowledge.

I'm glad you brought this up, Cassiodorus, because this is the most fundamental point. Otherwise, Tradition is just another point of view, another philosophy to debate, another infinite conversation without resolution. Our aim is always that intellectual conversion where we suddenly ``get it'' and can see the world in no other way.

\hfill

\texttt{Ash on 2013-08-23 at 03:07 said:}

If I might add, Cassiodorus, you mentioned faith in your post. This seems to be the key to overcoming this problem which we certainly share. Neither of us having attained gnosis personally, we are left in some sense to be shown the way by those who have. St. Thomas, whom you mentioned, believed that it made his Summa as ``piles of straw''. The scriptures say that we are saved ``by Grace through Faith''. Although we call Gnosis ``knowing'', and it is, I have always found it more useful to think of it as ``awareness''. As Buddha said, ``I am awake.''

We might have endless debates about physics or biology or social mores, which change with information. Gnosis, awareness of the union with the ground of all Being, cannot change. Those of us who have not attained it walk on having faith that it can be had, inspired by the example of the Saints, the Elect, who have attained it, and by the example of their works in the world. I've also been frequenting a certain Thomist site (incidentally, recommended by Cologero), and it certainly does change a lot of the questions to consider things from that perspective. As far as the commentor is concerned, I'd say it might be true that in the human state in particular, our first experience of knowledge is sensory. However, gnosis makes one aware of the fundamental ground of all Being. In other words, while we experience sensory knowledge first, Gnosis is knowledge of that which is metaphysically prior to the space-time material world from which we get it.

\hfill

\texttt{Cassiodorus on 2013-08-23 at 10:22 said:}

Cologero and Ash,

I appreciate your comments. They are helpful. In my last comment, I had remarked that many of my fellow Catholics have expressed the view that the Perennialist view is a modern day gnosticism. The reason being that the Christian cannot accept the doctrine of the Self and cannot interpret Creatio Ex Nihilo as Creatio Ex Deo. A Buddhist commenter that I shared this with said that there seems to be a confusion regarding the difference between non-duality and monism.

\hfill

\texttt{Cologero on 2013-08-23 at 12:46 said:}

Cassiodorus, perennialism has always been a part of Catholic theological thinking; are those fellows willing to jettison Plato and Aristotle? As long as they are stuck on the level of rational discourse, they cannot understand; that is a distinction Aquinas made. The problem is that few understand what that means any more. Your Buddhist friend is correct about that confusion. Dualism is decidedly un-Catholic, so non-dualism is better.

\hfill

\texttt{Gabe Ruth on 2013-08-23 at 15:39 said:}

Thank you, Cassiodorus, for the fruitful question.

I wonder what anyone here thinks of Zen and the Art of Motorcycle Maintenance by Robert Pirsig? I can't speak for his work after that book, and I won't claim that ZAMM is perfect. The author is beholden to the cult of the rebel and quite modern. Yet for all that, he sees the problems with modernity pretty well, and not in a utilitarian way. There's a passage where he casually ticks off various material concerns that progress had apparently taken care of for good, which sounds very odd from our current perspective.

He claims that the root of the problem is in ancient Greek philosophy, when Plato destroyed the Sophists and elevated dialectic and reason above rhetoric as the arbiter and revealer of Truth. Because we hear rhetoric and think politician, we're fully on board with that, and since all of Western philosophy is a footnote to Plato we all know the Sophists were no good. Pirsig believes that Plato (unintentionally) destroyed any concept of the good that could inspire humans by sterilizing it as eternal and unchanging, and that his victory left human knowledge at the mercy of categorizers such as his villian, Aristotle, because dialectic can be controlled by whoever makes up the categories that direct discursive thought. My first thought was that this is non-sense, we are ruled by empty rhetoric now, but as I thought about it two things occurred to me. In a more robust age, a people would not tolerate rulers whose actions and rhetoric were so far removed from each other. Rhetoric in the service of a noble soul has the ability to elevate man, instead of deceiving him. So excellence in rhetoric must not be isolated from excellence in every area, which is what Pirsig said the Sophists believed.

The second thing I think I came up with, but it is possible I just stole it from Dr. Charlton. Rhetoric is only a tool, even in the most skilled of hands, but dialectic and reason have the ability to become a tyrant and dominate the human mind, blocking out any other way of knowing.

\hfill

\texttt{Cassiodorus on 2013-08-24 at 00:47 said:}

I suspect that a traditional Catholic would say that the Traditionalists have ``stretched'' the Perennial Philosophy to include traditions that are totally alien to the West's Antique-Medieval synthesis. As I understand it, the Perennial Philosophy is supported by the profound unanimity found among the world's mystical traditions. But from this, most theists do not conclude that there is a transcendent unity of religions. That Thomist that I was speaking to said to me that the classical theist denies the notion that esoterism goes beyond theology, but actually enriches and deepens it.

\hfill

\texttt{Logres on 2013-08-24 at 09:26 said:}

I don't see how that the statement ``esotericism can't go beyond theology'' can really stand, without so much explanation and qualification as to render it useless for what they are using it for. Besides, they are taking credit for the riches and glories of esotericism a little belatedly -- the Church until this last century seems to have taken the point of view that outside the walls of the RCC, there was nothing. While true in one sense, and perhaps in a second sense of ``don't give a child a sword'', my perception (and maybe this is colored by being raised Protestant) is that it was meant in the literal sense. They were well-meaning, \& had the moral high ground in many ways, but it was ill-timed. If fathers aren't supposed children to wrath, I think the Church's insistence that esotericism was fully exteriorized in its current physical body was a mistake along these lines, and a symptom of the Kali Yuga itself. Why else was the Church Ichabod, glory departed, unless as a sign that judgement begins in the house of God? And yet, being able to appreciate and understand their position is also critical (at the very least, for the souls of our children), so I suppose dialogue with traditional Thomists is a very fruitful thing, and I salute your efforts, sir!

\hfill

\texttt{Logres on 2013-08-24 at 09:31 said:}

Was it someone on this Forum that first provided a link to Nicholas Gomez Davila? His meditations seem to be of the kind that Catholics could more fruitfully have pursued in this last century: 
\url{https://don-colacho.blogspot.com/}

I have to admit, however, am not very familiar with modern Thomism, as I only had one professor who ``followed'' them. He did have very correct opinions on the innate banality of basketball, however.

\hfill

\texttt{Cassiodorus on 2013-08-25 at 01:46 said:}

Logres,

``... ... ... I salute your efforts, sir!'' Thank you very much. Though probably undeserved. I came to the Christian faith by way of the writings of the traditionalist school. But, I must confess- I sometimes feel like a man without country in holding true to Perennialism and the tradition of my birth, Catholicism. ``Why so ?'' one might ask. Well, the vast majority of my fellow Catholics are not only indifferent to the Perennial Philosophy, but in many cases , downright hostile.

Perhaps, this kind of discussion is not appropriate here. Nevertheless, in the process of deepening my understanding of the Christian faith, i have so often come upon authoritative voices from within the Church that charge my esoteric dispositions with error and misunderstanding. Well, that's a hard pill to swallow on multiple levels.

My Thomist interlocutors seem to object to Perennialism on a few major counts:

* Inverting the relationship of metaphysics and revelation

* Putting the Trinity on the level of the ``relative Absolute''- not the Godhead

* The impossibility of the Second Person of the Trinity (the uncreated Logos) being plural as the created Logos- denying the uniqueness of Jesus Christ

* Pantheism

I suspect that a Traditionalist would say that the above points would be of no concern for the authentic esoterist and certainly not helpful for the ``one thing needful''. Probably so. Yet, they trouble me nonetheless.

\hfill

\texttt{Cologero on 2013-08-25 at 11:09 said:}

Cassiodorus, the main claim of Guenon is that there is a common metaphysical understanding behind the outer diversity of traditions. I believe that has been established. Otherwise, a Thomas Aquinas, inter alia, would have had nothing in common with Plato, Aristotle, the Neoplatonists, the Muslim philosopher, etc. To the contrary, he relied on them heavily. We have quoted writers from Augustine to Cardinal Manning on their understanding of a common tradition. So thus far there nothing objectionable. Actually, just the opposite, since contemporary theology is latching onto all sorts of erroneous philosophies as a support, as though theology were totally independent of any true metaphysics.

Now Thomism holds theology to be superior to to metaphysics, while Guenon reverses those roles. That is where the discussion belongs and may be of value. Nevertheless, Thomism, insofar as it relates to metaphysics and not revelation, can benefit by engaging Guenon's metaphysics. The accusation of pantheism is a calumny; Thomists should follow the angelic doctor is engaging a metaphysic that is not meant to be a rival, since, in the West, Thomism is the best expression of that metaphysic.

There is a deeper aspect. The goal of tradition is to achieve higher states; as such, the states are irrefutable and serve as the ``sensory data'' for further reflection. Now, those who achieve such states are seldom theologically sophisticated, so their descriptions, but not the fact, of such states may be insufficient from the theological perspective. There are many works that do the task of ``cleaning up'' the theology; nevertheless, the ``one thing necessary'' is achieve those states and not to be a well regarded theologian.

\hfill

\texttt{Michael on 2013-08-25 at 20:51 said:}

This has been a very helpful discussion as I am wrestling with some of the same questions that Cassiodorus has raised although I came to Catholicism first and then Perennialism recently.

My main issue with Perennialism as presented by Guenon is that it denies the uniqueness of Christ. Jesus claimed to be God. If his claim is true, then the Church he founded is unique and authoritative as well. It also has the uncomfortable corollary that subsequent `revelations' that contradict Christ's claims cannot be valid, even though there might be substantial truth in those traditions.

In other words, tradition is true and there are other traditional cultures, but Catholicism (and Orthodoxy) are true in a special way because there has been a divine intervention into history that changes the playing field. Before Christ, one should be part of his own traditional culture. After Christ, one is honor bound in accepting tradition of Christianity because it is true.

I am interested to hear how other Christians on this blog deal with these questions. I am not trying to convert any one here. I am just sharing my own difficulty and asking for guidance.

\hfill

\texttt{Cassiodorus on 2013-08-26 at 02:24 said:}

Cologero,

I'm grateful for your insightful and measured response.

I shared this exchange with a friend. He replied that the Perennialism, to him, is a failed attempt at assimilating classical theism to Vendatin non-dualism. There is a indeed, a Perennial Philosophy, he said, except that philosophy is not monistic. I answered by saying that I think nondualism is not equivalent to monism. When Perennialists say that all is one, they do not mean that all is of the same substance- rather that all is dependent and derivative of the One. To say that ``all is one'' is different than ``all is the same''. I think its also the case that non-dualism doesn't mean ``not two''. It would seem that ultimate reality is difficult to pin down.

\hfill

\texttt{Cologero on 2013-08-26 at 12:34 said:}

Cassiodorus, it is impossible to pin down in words; it can only be intuited.

\hfill

\texttt{Cassiodorus on 2013-08-27 at 17:49 said:}

Cologero,

I realize that you probably feel like this subject has already been discussed to exhaustion. But, if I may beg your indulgence, perhaps we can continue this just a little further? The Traditionalist William Stoddart has this to say in an essay entitled ``Mysticism'',

\emph{Metaphysically, the difference between exoterism and esoterism lies in how the final goal is envisaged: in exoterism (and in bhaktic esoterism), God is envisaged at the level of ``Being'' (the Creator and Judge): no matter how deep, how sublime the exoterist's fervor, Lord and worshipper remain distinct. In jnanic esoterism , on the other hand, God is envisaged on the level of ``Beyond-Being'' (the Divine Essence). At this level, it is perceived that Lord and worshipper (the latter known to be created in the image of the former) share a common essence, and this opens up the possibility of ultimate Divine union.}

As I understand it, Christianity denies outright that worshipper and God share a common essence. Classical theists would say, I think, that such a doctrine is \emph{precisely} what characterizes pantheism/Gnosticism. For the Christian, it's not a matter of spiritual temperament that distinguishes between jnanic and bhaktic mysticism, rather there simply cannot be such a thing as jnanic esoterism because, in essence, the human person is ontologically different from God. Nevertheless, we are told that ``the human person has nevertheless been created for communion with God, for Divinization: to become by grace or free gift what God is by nature, essentially.''

Is the difference between theosis and ``realization'' merely a semantic difference? Theistic writers seem to say no. Timothy Ware expresses it this way,

\emph{Nor does the human person, when it ``becomes god'', cease to be human : `We remain creatures while becoming god by grace, as Christ remained God when becoming man by the Incarnation.' The human being does not become God ``by nature'', but merely a ``created god'', a god ``by grace'' or ``by status''.}

I

\hfill

\texttt{Santiago on 2013-08-27 at 19:55 said:}

Cassiodorus,

If you haven't read them, you might want to look at ``Christianity and the Doctrine of Non-Dualism'' by an anonymous French Cistercian monk, and ``Christian Gnosis from St. Paul to Meister Eckhart'' by Wolfgang Smith.

Smith has much to say on the point which you identify as posing a difficulty for some Christians with regards to Perennialism, feeling as though they cannot ``interpret Creatio Ex Nihilo as Creatio Ex Deo''. I can give a clumsy summary of some of his points as I understand them, which I hope isn't out of place here:

A psalmist tells us ``God spoke once and for all and I have heard two things'' (Psalm 62:11). Meister Eckhart, in his Expositio Libri Genesis, interprets these two things as the emanation of the Persons of the Trinity and the creation of the world. So we see two things where God has only spoken once, which is the essence of creaturely misapprehension: ``every kind of existence outside intellect is creaturely''. The Gospel of John tells us that in the beginning was the Word, not in the beginning was Being. The Word implies subjectivity, interiority, depth-dimension -- In Eckhart's terms, intellect. Thus, the ``creaturely'' mistake is perceiving things as though they were outside, as mere objects. (I don't think Smith mentions it, but this corresponds to the Hermetic phrase ``to know is to be'', because knowing a thing righty is to experience its inwardness).

Thus, creation is not really an overflowing, but an inward flowing; not an ebullition but a ``bullition''. The nihilo out of which God creates would be His plenitude (white which contains all colour) as opposed to a complete black which lacks all colour.

Scripturally this nothingness might be found in Genesis 1:1. In the original Hebrew, the first words ``Bereshith bars Elohim'' (in the beginning, created Elohim), ``- reshith'' is the verb ``beginning'' and is identified by the Zohar as the second Sephiroth of the Cabala, Hokhmah. This verb isn't linked to a subject, so the Zohar interprets the subject as nothing: the unmanifest first Sephiroth, Kether Elyon. Elohim becomes the object of the phrase, which is rendered as ``in (by means of) Hokhmah, Kether created Elohim''. If Elohim corresponds to the lower Sephirah which constitute creation, these are all made by means of Hokhmah, so that creation itself is Divine stuff, our doors of perception are just too clogged up to see it.

Smith develops this in terms of the Christian Trinity. Vladimir Soloviev makes the same points, as I've understood him: everything must correspond to a conscious experience, and so reality takes place inside the ``broadest'' experiencer, God. But understanding this doesn't lead to our disappearing -- the world is still where it was, it's just apprehended properly.

Anyway, these are just thought-games which might be helpful, but as Cologero and you have pointed out, it all comes down to direct intuition. Getting our theology right is secondary.

\hfill

\texttt{Cassiodorus on 2013-08-28 at 18:10 said:}

Hi Santiago,

Thank for the recommendations. Turns out that I'm presently reading Wolfgang Smith's book and I own (but haven't read yet) ``Christianity and the Doctrine of Non-Duality''. For anyone interested in this subject, I'm also reading an excellent series of email exchanges between the Traditionalist Charles Upton and the Catholic (I think) Robert Bolton which can be found at Sacredweb under ``Faith Traditions and Ecumenism''.


\hfill

\end{sffamily}
\end{footnotesize}

